\subsection{Order-restriction DAGs}\label{sec:order-dags}

An \emph{order-restriction DAG} $D$ consists of a directed acyclic graph
on a finite set $V$, together with an open interval
$(\ell_x,r_x)\subseteq(0,1)$ for each vertex $x$. Its \emph{feasible
assignments} are the points in
\begin{equation}\label{eq:node-region}
K(D)=\{X\in\R^V:\ell_x<X_x<r_x\text{ for every }x,
\quad X_x<X_y\text{ for every arc }x\to y\}.
\end{equation}
Assignments with distinct coordinates determine relative orders of the
vertices. For example, the arcs $a\to b$ and $a\to c$, with all three
intervals equal to $(0,1)$, allow both $a<b<c$ and $a<c<b$.
Write $V(D)=\vol K(D)$, taking $V(D)=1$ when $V$ is empty.

We call the intervals \emph{canonical} when each is the full range of
its coordinate in feasible assignments. The routine
\texttt{Normalize}$(D)$ propagates lower bounds forward along arcs and
upper bounds backward. Two topological passes suffice, as shown in
Algorithm~\ref{alg:normalize}.

For a canonical $D$, the operation \texttt{Restrict}$(D,x,c)$ fixes
$X_x=c$ and removes its coordinate. For each arc $y\to x$, it sets
$r_y\gets\min(r_y,c)$, and for each arc $x\to z$, it sets
$\ell_z\gets\max(\ell_z,c)$. It then deletes $x$ and its incident arcs
and normalizes the remaining bounds. Algorithm~\ref{alg:restrict}
gives the procedure.

\begin{lemma}[Canonical intervals and coordinate restriction]\label{lem:interior}
Normalization preserves $K(D)$. This region has positive volume exactly
when all the resulting intervals have positive width. In that case the
intervals are canonical. For any $c\in(\ell_x,r_x)$, the feasible region
of \textnormal{\texttt{Restrict}}$(D,x,c)$ is exactly the section of
$K(D)$ with $X_x=c$, measured in the remaining coordinates, and has
positive volume.
\end{lemma}
Appendix~\ref{app:geometry} proves the lemma.

To measure the effect of restriction, suppose $V(D)>0$ and choose $X$
uniformly from $K(D)$. The density $f_x(c)$ of coordinate $X_x$ is the
volume of the section at $c$ divided by $V(D)$. Thus the lemma gives
\begin{equation}\label{eq:restriction-volume}
V(\operatorname{Restrict}(D,x,c))=V(D)f_x(c).
\end{equation}
This defines the density at every interior $c$ and gives
$\Prb(a<X_x<b)=\int_a^b f_x(t)\,dt$.
